\section{Conclusioni}
\label{sec:conclusioni}

Il progetto PomodorIA ha dimostrato, tramite un Proof of Concept interamente
software, la realizzabilità concettuale e pratica di un sistema di Edge AI
applicato alla Smart Agriculture, seguendo l'intero ciclo di vita che un
simile sistema affronterebbe in un contesto reale: dalla scelta e
validazione di un modello di Deep Learning appropriato al compito (Sezione
\ref{sec:modello-ml}), alla sua compressione per un ipotetico deployment su
hardware a risorse limitate (Sezione~\ref{sec:compressione}), fino
all'integrazione in un sistema completo di sensori, attuatori e un agente
decisionale razionale, formalizzato secondo il modello PEAS e collegato
esplicitamente ai riferimenti teorici richiamati nella
Sezione~\ref{sec:contesto}.

Il lavoro svolto non si è limitato all'implementazione della pipeline nel
suo funzionamento atteso, ma ha incluso un processo di verifica
sperimentale che ha portato all'identificazione e alla risoluzione di un
problema non evidente a un primo sguardo, il \emph{data leakage} nella
valutazione dell'accuratezza (Sezione~\ref{sec:moduli}).

I risultati sperimentali ottenuti (Sezione~\ref{sec:risultati}) evidenziano
come, per l'architettura convoluzionale relativamente compatta utilizzata in
questo progetto, la quantizzazione risulti la tecnica di compressione più
efficace, offrendo una riduzione della dimensione del modello di circa
quattro volte a fronte di una perdita di accuratezza trascurabile, mentre il
pruning strutturato, pur concettualmente rilevante, introduce in questo
caso specifico un compromesso meno favorevole.

Il progetto, pur trattandosi di una simulazione priva di hardware fisico,
è stato progettato secondo principi di ingegneria del software che ne
renderebbero possibile l'estensione verso un deployment reale con un
impatto limitato sulla logica applicativa, e apre a numerose direzioni di
sviluppo futuro discusse nella Sezione~\ref{sec:limiti}, dal
Federated Learning al Digital Twin, che confermano la validità del sistema
progettato come base per esplorazioni più ampie nell'ambito dell'Edge AI
per l'Internet of Things.
